\documentclass[10pt,a4paper]{amsart}
\usepackage[margin=19mm]{geometry}
\usepackage{amsmath,amssymb,mathtools}
\usepackage[scr=rsfs,cal=euler]{mathalfa}
\usepackage{xfrac}
\usepackage[unicode,hidelinks,bookmarks=false]{hyperref}
\newcommand{\ie}{i.e.\ }
\newcommand{\cf}{cf.\ }
\newcommand{\resp}{respectively\ }
\newcommand{\Subsection}{\subsection}
\providecommand{\nobreakdash}{\nobreak\hbox{-}}
\setlength{\emergencystretch}{2em}
\multlinegap0pt
\usepackage[shortlabels]{enumitem}
\usepackage{mathtools}
\usepackage{amssymb}
\usepackage{xcolor}
\usepackage{tikz}
\newtheorem{assumption}{Assumption}
\newtheorem{proposition}{Proposition}
\title{Size-Selective Threshold Harvesting under Nonlocal Crowding and Exogenous Recruitment}
\author{Jiguang Yu, Louis Shuo Wang, Ye Liang}
\date{Source: arXiv:2603.27368v2 (CC BY 4.0)}
\begin{document}
\maketitle

\noindent\emph{Source and licence.} Size-Selective Threshold Harvesting under Nonlocal Crowding and Exogenous Recruitment. Version arXiv:2603.27368v2 (CC BY 4.0), distributed by arXiv under the Creative Commons Attribution 4.0 International licence, http://creativecommons.org/licenses/by/4.0/, as recorded in the arXiv licence metadata for this exact version and shown on the abstract page https://arxiv.org/abs/2603.27368v2. Full licence notice: \texttt{licenses/arxiv-2603.27368-CC-BY-4.0.txt}.\par\smallskip

\noindent\textit{Editorial scope.} Counted unit: the article's proposition prop:necessary (source0.tex lines 683--724). The article's complete original proof of it is delivered with it (source0.tex lines 726--789, one \texttt{proof} environment of 64 lines). No mathematics is written, repaired or re-derived: the statement and the proof are the article's own lines, in the article's own notation, and no retained line is substituted or re-flowed.  Retained uncounted as the source-local context the proof needs: lines 287--328; lines 626--634; lines 640--646.  Nothing was omitted from any retained span, so the package carries no omission record.  No retained line contains a citation, so the wrapper carries no bibliography and no bibliography entry.  No retained line contains a figure environment, an image-inclusion command or drawing code, so no image is delivered and the package declares no asset dependency.  Editorial formatting only, in addition: an AMS \texttt{amsart} preamble that restates the article's own declarations and macros used by the retained lines, namely the environments \texttt{assumption}, \texttt{proposition} on separate counters, together with the standard packages those lines need.  The retained environments print in this document as Assumption 1, Proposition 1 in order of appearance, whereas the article prints the same items with the numbering of its own full document.  The complete native source of the article as distributed in the arXiv e-print for this exact version is bundled unchanged as \texttt{sources/197340/source0.tex}.
\begin{assumption}[Standing assumptions]\label{ass:standing}
\begin{enumerate}[label=(H\arabic*)]
\item\label{H:reg}
The functions $g(E,l)$ and $\mu(E,l)$ are continuous on
$[0,\infty)\times[l_0,l_m]$ and continuously differentiable with respect to $E$.

\item\label{H:gpos}
There exists a constant $g_{\min}>0$ such that
\[
g(E,l)\ge g_{\min}
\qquad \text{for all } (E,l)\in[0,\infty)\times[l_0,l_m].
\]

\item\label{H:mupos}
\[
\mu(E,l)\ge 0
\qquad \text{for all } (E,l)\in[0,\infty)\times[l_0,l_m].
\]

\item\label{H:mumono}
\[
\partial_E\mu(E,l)\ge 0
\qquad \text{for all } (E,l)\in[0,\infty)\times[l_0,l_m].
\]
That is, increased crowding does not decrease natural mortality.

\item\label{H:chi}
The weighting kernel $\chi(l)$ is continuous on $[l_0,l_m]$, satisfies
$\chi(l)\ge 0$ for all $l\in[l_0,l_m]$, and is not identically zero.

\item\label{H:uadm}
The harvesting control $u$ is measurable and satisfies
\[
0\le u(t,l)\le u_{\max}
\qquad \text{for all } (t,l)\in[0,\infty)\times[l_0,l_m].
\]

\item\label{H:ic}
The initial condition satisfies $\phi\in L^1([l_0,l_m])$ and
$\phi(l)\ge 0$ for a.e.\ $l\in[l_0,l_m]$.
\end{enumerate}
\end{assumption}

\begin{align}\label{eq:state_pmp}
\begin{aligned}
&\frac{\partial x(t,l)}{\partial t}
+\frac{\partial}{\partial l}\Big(g(E(t),l)\,x(t,l)\Big) \\[4pt]
={}&
-\mu(E(t),l)\,x(t,l)-u(t,l)\,x(t,l),
\qquad t>0,\; l\in(l_0,l_m),
\end{aligned}
\end{align}

\begin{equation}\label{eq:objective_J}
J[u]
\coloneqq
\int_{0}^{\infty} e^{-rt}
\int_{l_0}^{l_m} c(l)\,u(t,l)\,x(t,l)\,dl\,dt,
\qquad r>0,
\end{equation}

\begin{proposition}[Formal first-order necessary conditions]
\label{prop:necessary}
Assume that the coefficients, controls, and state variables satisfy
Assumption~\ref{ass:standing} and are sufficiently smooth for integration by parts to
be justified formally. Let $(\hat u, \hat x, \hat E)$ be an optimal solution to the problem
\eqref{eq:objective_J} and \eqref{eq:state_pmp} subject to
$g(\hat E(t),l_0)\,\hat x(t,l_0)=p(t)$. 
Then, formally, there exists an adjoint variable
$\lambda:[0,\infty)\times[l_0,l_m]\to\mathbb{R}$ satisfying:
\begin{align}\label{eq:adjoint}
\partial_t \lambda(t,l) &+ g(\hat E(t),l)\,\partial_l \lambda(t,l) \notag\\
&= \bigl(r + \mu(\hat E(t),l) + \hat u(t,l)\bigr)\lambda(t,l) - c(l)\,\hat u(t,l)
\notag \\
&\quad + \chi(l)\!\int_{l_0}^{l_m} \hat x(t,s) \Big(
   \lambda(t,s)\,\partial_E\mu(\hat E(t),s)
   - \partial_s\lambda(t,s)\,\partial_Eg(\hat E(t),s) \Big)\,ds.
\end{align}

The terminal-size and transversality conditions are:
\begin{align}
\lambda(t,l_m) &= 0 \qquad \text{for } t \ge 0,
   \label{eq:adjoint_boundary} \\
\lim_{T\to\infty} \int_{l_0}^{l_m} e^{-rT}\lambda(T,l)\,\delta x(T,l)\,dl &= 0
   \qquad \text{for every admissible variation } \delta x.
   \label{eq:transversality}
\end{align}

Define the switching function
\begin{equation}\label{eq:switching}
\Phi(t,l) \coloneqq \hat x(t,l)\bigl( c(l) - \lambda(t,l) \bigr).
\end{equation}
Assuming $\hat x(t,l)>0$ for a.e.\ $(t,l)$, the optimal harvesting profile satisfies the
bang--bang rule:
\begin{equation}\label{eq:opt_control}
\hat u(t,l)=
\begin{cases}
0, & c(l)<\lambda(t,l),\\[4pt]
u_{\max}, & c(l)>\lambda(t,l),\\[4pt]
\text{any value in }[0,u_{\max}], & c(l)=\lambda(t,l).
\end{cases}
\end{equation}
\end{proposition}

\begin{proof}
Define the Lagrangian by adjoining \eqref{eq:state_pmp} to the objective
\eqref{eq:objective_J} via the multiplier $\lambda(t,l)$:
\begin{align*}
L(x, u, \lambda) 
&= \int_0^\infty e^{-rt} \int_{l_0}^{l_m} c(l)u(t,l)x(t,l)\,dl\,dt \\
&\quad - \int_0^\infty e^{-rt} \int_{l_0}^{l_m} \lambda(t,l) \Big[ \partial_t x(t,l) + \partial_l\bigl(g(E(t),l)x(t,l)\bigr) \\
&\hspace{4.5cm} + \bigl(\mu(E(t),l) + u(t,l)\bigr)x(t,l) \Big]\,dl\,dt.
\end{align*}
Consider admissible perturbations $x \to x + \delta x$ and $u \to u + \delta u$, which induce a variation in the nonlocal environment variable $E \to E + \delta E$ where
\[
\delta E(t) = \int_{l_0}^{l_m} \chi(l)\delta x(t,l)\,dl.
\]
The variation of the Lagrangian with respect to the state, $\delta_x L$, must vanish at the optimum:
\begin{align*}
\delta_x L &= \int_0^\infty e^{-rt} \int_{l_0}^{l_m} \delta x \big(c(l)u\big)\,dl\,dt \\
&\quad - \int_0^\infty e^{-rt} \int_{l_0}^{l_m} \lambda \Big[ \partial_t \delta x + \partial_l(g \delta x + x \partial_Eg \delta E) + (\mu + u)\delta x + x \partial_E\mu \delta E \Big]\,dl\,dt = 0.
\end{align*}
We evaluate the integral terms utilizing integration by parts. For the time derivative, assuming $\delta x(0,l) = 0$ and invoking the transversality condition \eqref{eq:transversality} as $t\to\infty$, we get
\[
-\int_0^\infty e^{-rt} \lambda(t,l) \partial_t \delta x\,dt = \int_0^\infty e^{-rt} \bigl(\partial_t \lambda - r\lambda\bigr)\delta x\,dt.
\]
For the spatial derivatives, integration by parts yields:
\begin{align*}
-\int_{l_0}^{l_m} \lambda \partial_l\big(g \delta x + x \partial_Eg \delta E\big)\,dl 
&= - \Big[ \lambda \big(g \delta x + x \partial_Eg \delta E\big) \Big]_{l_0}^{l_m} 
+ \int_{l_0}^{l_m} \partial_l \lambda \big(g \delta x + x \partial_Eg \delta E\big)\,dl.
\end{align*}
Under the terminal-size condition $\lambda(t,l_m) = 0$, the boundary term at $l_m$ vanishes. 
At the recruitment boundary $l_0$, the admissible variations must preserve the inflow flux condition
\[
g(E(t),l_0)x(t,l_0)=p(t),
\]
with $p(t)$ exogenous. Therefore,
\[
\delta\!\big(g(E(t),l_0)x(t,l_0)\big)=0,
\]
that is,
\[
g(E(t),l_0)\,\delta x(t,l_0)+x(t,l_0)\,\partial_E g(E(t),l_0)\,\delta E(t)=0.
\]
Hence the boundary contribution at $l_0$ vanishes:
\[
-\Big[\lambda(t,l)\big(g\,\delta x+x\,\partial_E g\,\delta E\big)\Big]_{l=l_0}=0.
\]

Isolating all terms proportional to $\delta E(t)$ from the spatial and mortality variations, we define the environment coupling factor $C(t)$:
\[
C(t) = \int_{l_0}^{l_m} x(t,s)\Big[ \partial_s\lambda(t,s)\partial_Eg(E(t),s) - \lambda(t,s)\partial_E\mu(E(t),s) \Big]\,ds.
\]
Substituting $\delta E(t) = \int_{l_0}^{l_m} \chi(l)\delta x(t,l)\,dl$ back into the variational equation, this contributes $\chi(l)C(t)$ inside the spatial integral with respect to $\delta x(t,l)$.

Collecting all coefficients for $\delta x(t,l)$ inside the double integral gives:
\begin{align*}
\delta_x L &= \int_0^\infty e^{-rt} \int_{l_0}^{l_m} \delta x(t,l) \Bigg[ c(l)u + \partial_t \lambda - r\lambda + g\partial_l\lambda - \lambda(\mu + u) + \chi(l)C(t) \Bigg]\,dl\,dt = 0.
\end{align*}
Since this must hold for any arbitrary variation $\delta x$, the expression in the square brackets must vanish a.e. Rearranging this yields the adjoint equation \eqref{eq:adjoint}.

Finally, taking the variation with respect to the control $u$, we obtain:
\[
\delta_u L = \int_0^\infty e^{-rt} \int_{l_0}^{l_m} x(t,l)\big( c(l) - \lambda(t,l) \big)\delta u(t,l)\,dl\,dt.
\]
This defines the switching function $\frac{\partial L}{\partial u}$, dictating the optimal bang-bang policy described in \eqref{eq:opt_control}.
\end{proof}

\end{document}
